\documentclass[11pt]{article}

\usepackage[margin=1in]{geometry}
\usepackage{amsmath,amssymb,amsfonts,amsthm}
\usepackage{array}
\usepackage{bm}
\usepackage{multirow}
\usepackage{multicol}
\usepackage{nicefrac}
\usepackage{paralist}
\usepackage{enumitem}
\usepackage{verbatim}
\usepackage{thm-restate}
\usepackage[normalem]{ulem}
\usepackage[compact]{titlesec}
\usepackage{fancybox}
\newenvironment{fminipage}%
  {\begin{Sbox}\begin{minipage}}%
  {\end{minipage}\end{Sbox}\fbox{\TheSbox}}

\newenvironment{algbox}[0]{%
\vskip 0.2in
\noindent
\begin{fminipage}{6.3in}
}{%
\end{fminipage}
\vskip 0.2in
}

\usepackage{xcolor}
\usepackage{cite}
\usepackage{url}
\usepackage{xspace}
\usepackage[mathscr]{euscript}
\usepackage{nameref}
\usepackage[
  linktocpage=true,
  pagebackref=false,
  colorlinks=true,
  linkcolor=blue,
  citecolor=blue,
  urlcolor=blue,
  bookmarks,
  bookmarksopen,
  bookmarksnumbered
]{hyperref}
\usepackage[noabbrev,nameinlink]{cleveref}

% Make the font fallbacks used by Computer Modern explicit and silent.
\DeclareFontShape{OT1}{cmr}{m}{scit}{<->ssub*cmr/m/sc}{}
\DeclareFontShape{OT1}{cmr}{bx}{sc}{<->ssub*cmr/bx/n}{}

\newtheorem{theorem}{Theorem}[section]
\newtheorem{example}{Example}[section]
\newtheorem{lemma}{Lemma}[section]
\newtheorem{proposition}[lemma]{Proposition}
\newtheorem{corollary}[lemma]{Corollary}
\newtheorem{claim}[lemma]{Claim}
\newtheorem{fact}[lemma]{Fact}
\newtheorem{conjecture}[lemma]{Conjecture}
\newtheorem{definition}[lemma]{Definition}
\newtheorem{problem}{Problem}
\newtheorem{remark}[lemma]{Remark}
\newtheorem{observation}[lemma]{Observation}

\newtheorem*{theorem*}{Theorem}
\newtheorem*{claim*}{Claim}
\newtheorem*{proposition*}{Proposition}
\newtheorem*{lemma*}{Lemma}
\newtheorem*{problem*}{Problem}

\allowdisplaybreaks
\setlength{\parskip}{3pt}
\setlength{\emergencystretch}{2em}

\title{Universal Contextual Composition Is Strictly Weaker than
Neighbor-Preserving Differential Obliviousness}
\author{}
\date{September 1, 2026}

\begin{document}

\pagenumbering{roman}

\maketitle

\begin{abstract}
For a randomized mechanism $M$, let $P_x$ be the law of the first-stage execution
$\operatorname{Exec}_M(x)=(V,Y)$ on a finite or countable space
$\Omega=\mathcal V\times\mathcal Y$, where the input and output spaces have
symmetric neighbor relations.  Write $aRb$ when two executions have the same
view and neighboring outputs, and let $R(S)=\{b:\text{some }a\in S\text{
 satisfies }aRb\}$.  For $\eta\geq0$ and $\gamma\in[0,1]$, an
$(\eta,\gamma)$-admissible context is a Markov kernel $K$ satisfying
$K(a,A)\leq e^\eta K(b,A)+\gamma$ for every measurable event $A$ whenever
$aRb$.  We define universal
contextual composability by quantifying over all such contexts.  For
$\varepsilon\geq0$ and $\delta\in[0,1]$, if the
$(\varepsilon,\delta)$-NPDO inequalities
$P_x(S)\leq e^\varepsilon P_{x'}(R(S))+\delta$ hold for every ordered pair
$x\sim_{\mathcal X}x'$ and every $S\subseteq\Omega$, then, for every
$\eta\geq0$, $\gamma\in[0,1]$, every $(\eta,\gamma)$-admissible $K$, and
every contextual event $A$,
\[
 (P_xK)(A)
 \leq e^{\varepsilon+\eta}(P_{x'}K)(A)
      +\delta+\gamma-\delta\gamma
 \qquad (x\sim_{\mathcal X}x').
\]
The converse fails with $\varepsilon=\delta=0$.  A finite mechanism with a
fixed view exchanges masses $1/3$ and $2/3$ between the endpoints of a
three-point output path.  It is $(0,0)$-universally contextually composable;
in fact,
for every $\eta\geq0$, $\gamma\in[0,1]$, $(\eta,\gamma)$-admissible kernel,
and event, its two nontrivial contextual comparisons have additive loss at most
\[
 \frac{\max\{0,2+e^\eta-e^{2\eta}\}}{3}\,\gamma
 \leq \frac{2\gamma}{3}.
\]
Nevertheless, it violates $(0,0)$-NPDO.  The mechanism is implemented in
$O(1)$ time using one exact uniform sample from three outcomes and outputs an
element of the three-point path with a fixed execution view.
\end{abstract}

\noindent\textbf{Note.} This paper was fully AI generated through a harness created by Richard Peng that discovers open problems and runs them through an automated research pipeline.\par

\clearpage
\tableofcontents
\clearpage

\pagenumbering{arabic}

\section{Introduction}
\label{sec:introduction}

Composability is particularly delicate when privacy protects what an
execution reveals but not the intermediate value that it computes.  Classical
oblivious computation seeks to hide data-dependent access patterns, whereas
\emph{differential obliviousness} (DO) asks the observable access pattern to
satisfy differential privacy on neighboring inputs
\cite{GoldreichOstrovsky1996ORAM,DworkEtAl2006DP,ChanEtAl2022FoundationsDO}.
This relaxation supports efficient algorithms for tasks such as database joins
and relational operators \cite{ChuEtAl2021DOJoins,QinEtAl2022Adore}.  In a
pipeline, however, a first algorithm may expose a view $V$ while passing an
unobserved output $Y$ to a second algorithm.  Privacy of $V$ alone does not
ensure that neighboring first-stage inputs produce neighboring values of $Y$,
so the second algorithm's privacy guarantee cannot simply be invoked.

Zhou, Shi, Chan, and Maimon introduced neighbor-preserving differential
obliviousness (NPDO) to control precisely this interface
\cite{ZhouEtAl2023NPDO}.  Their main composition theorem shows, on finite or
countable execution spaces, that an NPDO first stage composes additively with a
DO second stage.  The paper also characterizes NPDO by a supported fractional
matching between the two first-stage execution laws.  Subsequent work develops
advanced composition and divergence-based variants of NPDO
\cite{ZhouEtAl2024AdvancedNPDO}.  These results establish NPDO as a sufficient
interface.  They motivate a converse question: if a first stage remains private
after every admissible downstream context, must it satisfy NPDO?

In the finite-or-countable statistical model, we show that NPDO suffices for a
strong guarantee against every common downstream context but is not necessary:
the converse already fails with zero first-stage error on a three-point output
path.  To state the result, let the input and output spaces carry symmetric
neighbor relations, let
$\operatorname{Exec}_M(x)=(V,Y)$ have law $P_x$ on a finite or countable
space $\Omega=\mathcal V\times\mathcal Y$, and write
\[
 (v,y)R(v',y')
 \quad\Longleftrightarrow\quad
 v=v'\ \text{ and }\ y\sim_{\mathcal Y}y'.
\]
For $S\subseteq\Omega$, let $R(S)$ be the set of executions related to some
element of $S$.  The $(\varepsilon,\delta)$-NPDO condition is the setwise
requirement
\[
 P_x(S)\leq e^\varepsilon P_{x'}(R(S))+\delta
 \qquad(x\sim_{\mathcal X}x')
\]
for every $S\subseteq\Omega$.  An $(\eta,\gamma)$-admissible context is a
Markov kernel
$K:\Omega\rightsquigarrow\mathcal W$ into an arbitrary measurable observation
space such that
\[
 K(a,A)\leq e^\eta K(b,A)+\gamma
 \qquad(aRb)
\]
for every measurable event $A$.  We call $M$
$(\varepsilon,\delta)$-\emph{universally contextually composable} when, for
every $\eta\geq0$ and $\gamma\in[0,1]$, every such common kernel applied in
both neighboring worlds satisfies
\[
 (P_xK)(A)
 \leq e^{\varepsilon+\eta}(P_{x'}K)(A)+\delta+\gamma
 \qquad(x\sim_{\mathcal X}x').
\]
Our main theorem proves sufficiency and separates it from necessity.

\begin{theorem}[Sufficiency and strict separation]
\label{thm:main}
Fix $\varepsilon\geq0$ and $\delta\in[0,1]$.  If $M$ is
$(\varepsilon,\delta)$-NPDO on a finite or countable execution space, then for
every $\eta\geq0$, $\gamma\in[0,1]$, every
$(\eta,\gamma)$-admissible Markov context $K$, and every measurable contextual
event $A$,
\[
 (P_xK)(A)
 \leq e^{\varepsilon+\eta}(P_{x'}K)(A)
      +\delta+\gamma-\delta\gamma
\]
for every ordered pair $x\sim_{\mathcal X}x'$.  Hence $M$ is
$(\varepsilon,\delta)$-universally contextually composable.

The converse fails with zero first-stage error.  Let
$\mathcal X=\{x_0,x_1\}$ with every pair related, let
$\mathcal Y=\{0,1,2\}$ with $i\sim_{\mathcal Y}j$ exactly when
$|i-j|\leq1$, and let the execution view always equal $v$.  On input $x_i$,
the algorithm \hyperref[alg:three-point]{\textsc{ThreePoint}} outputs a vertex
$y\in\mathcal Y$ and realizes the execution laws
\[
 \begin{array}{c|ccc}
  & (v,0)&(v,1)&(v,2)\\ \hline
  P_{x_0}&1/3&0&2/3\\
  P_{x_1}&2/3&0&1/3.
 \end{array}
\]
It runs in $O(1)$ time assuming unit-cost exact uniform sampling from three
outcomes.  The resulting mechanism is $(0,0)$-universally contextually
composable but is not $(0,0)$-NPDO.  More strongly, for every $\eta\geq0$,
$\gamma\in[0,1]$, every
$(\eta,\gamma)$-admissible $K$, and every measurable $A$, define
\[
 c_{\eta,\gamma}
 :=\frac{\max\{0,2+e^\eta-e^{2\eta}\}}{3}\,\gamma
 \leq \frac{2\gamma}{3}.
\]
Then
\[
 (P_{x_0}K)(A)\leq e^\eta(P_{x_1}K)(A)+c_{\eta,\gamma}
 \quad\text{and}\quad
 (P_{x_1}K)(A)\leq e^\eta(P_{x_0}K)(A)+c_{\eta,\gamma}.
\]
\end{theorem}

This common-kernel definition allows the context to depend on the full
execution pair $(V,Y)$ and therefore includes the usual output-fed serial
composition as a special case.  We use the phrase ``universal contextual
composability'' only for this definition.  Thus, under the common-context
definition above, universal contextual composability is strictly weaker than
NPDO already for a finite output graph and
$(\varepsilon,\delta)=(0,0)$.

\paragraph{Relation to prior work.}
The direct NPDO results concern the same hidden-output interface as the present
theorem: they prove sufficient composition theorems, while our result asks what
is forced by quantification over all common downstream contexts.  Structured
forms of the interface also occur in guarantees for threshold binning,
local-randomizer-to-DO-shuffler pipelines, and distance-preserving database
operators
\cite{ChanEtAl2022FoundationsDO,GordonEtAl2022PrivacyBlanket,ZhouShi2022DOShuffle,QinEtAl2022Adore}.
Those are restricted, mechanism-specific composition results rather than
characterizations of arbitrary first-stage execution laws.

The setwise NPDO inequality also has a nearby, non-DO-specific interpretation.
After aligning one-sided and symmetric conventions, it is the discrete form of
a witness-free approximate relational lifting
\cite{Sato2016ApproximateRHL}; on discrete spaces, approximate Strassen
theorems turn this inequality into a supported star-lifting or fractional
matching \cite{BartheEtAl2019StarLiftings,ZhouEtAl2023NPDO}.  Universal
relational-testing formulations in the codensity-lifting literature quantify
over a \emph{pair}
of related contexts \cite{KatsumataEtAl2018Codensity,AguirreEtAl2021Adversarial}.
Our property instead uses one common kernel in both worlds.  It is this
diagonal restriction, together with a nontransitive output adjacency relation,
that leaves room for the converse to fail.

\paragraph{Proof ideas.}
For sufficiency, the NPDO neighborhood inequalities yield a fractional
$R$-subcoupling that matches all but $\delta$ of $P_x$ and whose right marginal
is bounded by $e^\varepsilon P_{x'}$.  Transporting the function
$a\mapsto K(a,A)$ through this subcoupling keeps the exact unmatched mass
rather than replacing it immediately by $\delta$.  The contextual error is
then charged only on matched mass, giving the sharper bound
$\delta+\gamma-\delta\gamma$ without exponential amplification.  For strictness, the
three-point path has no mass at its middle vertex.  The endpoint laws fail the
one-hop neighborhood comparison required by NPDO, while admissibility still
relates the two endpoints through the middle vertex; averaging those two-hop
constraints according to the displayed laws gives the sharper contextual
bound in \Cref{thm:main}.

\paragraph{Organization.}
\Cref{sec:preliminaries} formalizes execution laws, common contexts, NPDO, and
fractional subcouplings.  \Cref{sec:overview} explains the
matching-and-transport argument and the three-point separation.
\Cref{sec:main-proof} proves both parts of \Cref{thm:main}.

\section{Preliminaries}
\label{sec:preliminaries}

\paragraph{Execution laws and relational neighborhoods.}
Let $(\mathcal X,\sim_{\mathcal X})$ and
$(\mathcal Y,\sim_{\mathcal Y})$ be sets with symmetric neighbor relations.
For a randomized mechanism $M$, write
$\operatorname{Exec}_M(x)=(V,Y)$ for its execution record on input $x$, where
$V\in\mathcal V$ is the observable view and $Y\in\mathcal Y$ is the output.
Throughout, the execution space
$\Omega:=\mathcal V\times\mathcal Y$ is finite or countable and has the
discrete $\sigma$-algebra.  Its execution law is
\[
 P_x(a):=\Pr[\operatorname{Exec}_M(x)=a].
\]
For any probability mass function $p$ on $\Omega$ and $S\subseteq\Omega$,
write $p(S):=\sum_{a\in S}p(a)$.
For $a=(v,y)$ and $b=(v',y')$, define the symmetric execution relation
\begin{equation}
 aRb
 \quad\Longleftrightarrow\quad
 v=v'\ \text{ and }\ y\sim_{\mathcal Y}y',
 \label{eq:execution-relation}
\end{equation}
and, for $S\subseteq\Omega$, its relational neighborhood
\begin{equation}
 R(S):=\{b\in\Omega:\text{there exists }a\in S\text{ with }aRb\}.
 \label{eq:relational-neighborhood}
\end{equation}
All privacy comparisons below are required for each ordered neighboring pair
$x\sim_{\mathcal X}x'$.

\paragraph{Contexts.}
Let $(\mathcal W,\mathcal F)$ be an arbitrary measurable observation space.  A
Markov kernel $K:\Omega\rightsquigarrow\mathcal W$ assigns a probability
measure $K(a,\cdot)$ to every $a\in\Omega$.  For a law $P$ on $\Omega$, the
contextual law $PK$ is
\begin{equation}
 (PK)(A):=\sum_{a\in\Omega}P(a)K(a,A),
 \qquad A\in\mathcal F.
 \label{eq:contextual-law}
\end{equation}

\begin{definition}[NPDO and universal contextual composability]
\label{def:npdo-ucc}
Fix $\varepsilon\geq0$ and $\delta\in[0,1]$.
The mechanism $M$ is $(\varepsilon,\delta)$-NPDO if, for every
$x\sim_{\mathcal X}x'$ and every $S\subseteq\Omega$,
\begin{equation}
 P_x(S)\leq e^\varepsilon P_{x'}(R(S))+\delta.
 \label{eq:npdo}
\end{equation}
For $\eta\geq0$ and $\gamma\in[0,1]$, a kernel
$K:\Omega\rightsquigarrow\mathcal W$ is
$(\eta,\gamma)$-admissible if, for every $aRb$ and $A\in\mathcal F$,
\begin{equation}
 K(a,A)\leq e^\eta K(b,A)+\gamma.
 \label{eq:admissible}
\end{equation}
The mechanism $M$ is $(\varepsilon,\delta)$-universally contextually
composable if, for every $\eta\geq0$, $\gamma\in[0,1]$, every observation
space, every $(\eta,\gamma)$-admissible $K$, every
$x\sim_{\mathcal X}x'$, and every $A\in\mathcal F$,
\begin{equation}
 (P_xK)(A)
 \leq e^{\varepsilon+\eta}(P_{x'}K)(A)+\delta+\gamma.
 \label{eq:ucc}
\end{equation}
\end{definition}

The NPDO clause in \Cref{def:npdo-ucc} is the discrete setwise formulation of
Zhou, Shi, Chan, and Maimon~\cite{ZhouEtAl2023NPDO}.  The
universal-composability clause fixes the meaning of that term throughout this paper.
It includes the usual output-fed serial context.  Indeed, if a downstream
algorithm takes $y$ as input and has view kernel $L_y$ on a measurable space
$\mathcal Z$, retain the first view and use observations in
$\mathcal V\times\mathcal Z$.  For an event $A$ and a fixed $v$, set
$A_v:=\{z:(v,z)\in A\}$ and define
\[
 K((v,y),A):=L_y(A_v).
\]
Whenever the kernels $L_y$ and $L_{y'}$ satisfy the downstream privacy
inequality for $y\sim_{\mathcal Y}y'$, this $K$ is admissible by
\Cref{eq:execution-relation,eq:admissible}.

\paragraph{Fractional subcouplings.}
Let $p$ and $q$ be probability mass functions on a finite or countable set
$\Omega$, let $R\subseteq\Omega\times\Omega$, and let $c\geq1$.  A
$c$-dominated fractional $R$-subcoupling of $p$ and $q$ is a function
$w:\Omega\times\Omega\to\mathbb R_{\geq0}$ satisfying
\begin{align}
 w(a,b)>0&\quad\Longrightarrow\quad aRb,
 \label{eq:w-support}\\
 \sum_b w(a,b)&\leq p(a) &&\text{for every }a\in\Omega,
 \label{eq:w-row}\\
 \sum_a w(a,b)&\leq c q(b) &&\text{for every }b\in\Omega.
 \label{eq:w-column}
\end{align}
Its mass is $\operatorname{mass}(w):=\sum_{a,b}w(a,b)$.  The row bound gives
$\operatorname{mass}(w)\leq1$; we call
$1-\operatorname{mass}(w)$ its unmatched mass.

\begin{fact}[Countable approximation and finite flows]
\label{fact:basic-tools}
We use the following standard facts.  If $h:D\to\mathbb R_{\geq0}$ is
summable on a countable set, then for every $\rho>0$ there is a finite
$F\subseteq D$ with $\sum_{z\notin F}h(z)<\rho$.  If $D$ is countable, every
sequence in $[0,1]^D$ has a pointwise-convergent subsequence.  Finally, in a
finite network with nonnegative capacities, the maximum flow value equals the
minimum capacity of a source--sink cut.
\end{fact}

\section{Overview}
\label{sec:overview}

The proof of \Cref{thm:main} has two parts.  For sufficiency, we turn the
setwise NPDO inequalities into a fractional matching between neighboring
execution laws and then transport an arbitrary contextual event through that
matching.  The central difficulty is the error accounting: NPDO controls
indicators of sets, while a context tests the law against the $[0,1]$-valued
function $a\mapsto K(a,A)$, and the desired conclusion must have additive
loss at most $\delta+\gamma-\delta\gamma$ without multiplying either error by
a privacy factor.  For strictness, we construct a finite endpoint-exchange mechanism on
a three-point path.  Its laws fail the one-hop neighborhood comparison
required by NPDO, although every admissible context can compare them through
the two edges of the path.

\paragraph{The fractional matching lift.}
Fix an ordered pair $x\sim_{\mathcal X}x'$ and apply the NPDO inequality with
$p=P_x$, $q=P_{x'}$, and $c=e^\varepsilon$.  The countable fractional
matching lift, \Cref{lem:matching}, produces weights $w(a,b)$ satisfying
\[
 \begin{gathered}
 w(a,b)>0\ \Longrightarrow\ aRb,
 \qquad \sum_b w(a,b)\leq P_x(a),
 \qquad \sum_a w(a,b)\leq e^\varepsilon P_{x'}(b),\\
 \sum_{a,b}w(a,b)\geq1-\delta.
 \end{gathered}
\]
Thus $w$ couples all but at most $\delta$ of the first law to related
executions, while allowing the second law's capacity to be enlarged by the
multiplicative factor $e^\varepsilon$.  On a finite space, these weights are
the flow in the bipartite graph induced by $R$: the NPDO set inequality is
precisely the cut condition needed by max-flow/min-cut.  For countable
$\Omega$, we build flows on finite truncations with mass $1-\delta-o(1)$ and
take a coordinatewise limit.  A tightness argument prevents matched mass from
escaping to infinity, which is the extra point needed beyond the finite
fractional Hall argument.

\paragraph{Contextual transport.}
For an $(\eta,\gamma)$-admissible context, fix an event $A$, write
$k(a)=K(a,A)$, and let $m=\sum_{a,b}w(a,b)$.  The four properties above give
the complete transport estimate
\begin{align*}
 (P_xK)(A)
 &\leq \sum_{a,b}w(a,b)k(a)+1-m\\
 &\leq e^\eta\sum_{a,b}w(a,b)k(b)+\gamma m+1-m\\
 &\leq e^{\varepsilon+\eta}(P_{x'}K)(A)
       +\delta+\gamma-\delta\gamma.
\end{align*}
The support of $w$ permits the pointwise admissibility comparison, its column
bound supplies $e^\varepsilon$, and $m\geq1-\delta$ gives
$\gamma m+1-m\leq\delta+\gamma-\delta\gamma$.  Thus contextual error is paid
only on matched mass, while the unmatched mass is not multiplied by $e^\eta$.
Since the context,
event, and ordered neighboring inputs were arbitrary, this proves the positive
direction through \Cref{lem:transport}.

\paragraph{The separating mechanism.}
For the converse, let the output relation be the path
$0\sim_{\mathcal Y}1\sim_{\mathcal Y}2$ and let the execution view always be
$v$.  The two neighboring inputs have laws
\[
 \begin{aligned}
 \bigl(P_{x_0}((v,0)),P_{x_0}((v,1)),P_{x_0}((v,2))\bigr)
   &=(1/3,0,2/3),\\
 \bigl(P_{x_1}((v,0)),P_{x_1}((v,1)),P_{x_1}((v,2))\bigr)
   &=(2/3,0,1/3).
 \end{aligned}
\]
The \hyperref[alg:three-point]{\textsc{ThreePoint}} algorithm realizes these
laws exactly: on input $x_i$ it takes one uniform ternary sample, selects the
corresponding endpoint, and writes it using a fixed constant-length access
schedule whose view is $v$.  Its output is the selected path vertex, its joint
execution law is the displayed row, and it runs in $O(1)$ time under unit-cost
exact ternary sampling.  The zero mass at the middle vertex makes the NPDO
failure immediate.  For $S=\{(v,2)\}$,
\[
 R(S)=\{(v,1),(v,2)\},
 \qquad
 P_{x_0}(S)=\frac23>\frac13=P_{x_1}(R(S)).
\]

\paragraph{Why every admissible context still composes.}
The unused middle vertex nevertheless constrains a context at both endpoints.
For an arbitrary contextual event, set $t=e^\eta$ and
\[
 a=K((v,0),A),\qquad b=K((v,1),A),\qquad c=K((v,2),A).
\]
Chaining admissibility along the two edges gives
\[
 c\leq tb+\gamma\leq t^2a+(t+1)\gamma,
 \qquad
 a\leq tb+\gamma\leq t^2c+(t+1)\gamma.
\]
A bare two-hop comparison would lose a second factor of $t$.  The asymmetric
endpoint masses are chosen so that their averaging removes this loss.  Indeed,
for the two contextual probabilities
$\mu=(a+2c)/3$ and $\nu=(2a+c)/3$, the key identity is
\[
 3(\mu-t\nu)=(1-2t)a+(2-t)c.
\]
Combining this identity with the appropriate two-hop bound, and then
exchanging $a$ and $c$ for the reverse comparison, yields
\[
 \mu\leq t\nu+
 \frac{\max\{0,2+t-t^2\}}{3}\,\gamma.
\]
The coefficient is at most $2/3$, so \Cref{lem:separation} proves the stronger
additive guarantee
$\max\{0,2+e^\eta-e^{2\eta}\}\gamma/3\leq2\gamma/3$, and hence
$(0,0)$-universal contextual composability.  Together with the one-hop NPDO
violation, this completes the strict separation.

\section{Proof of the Main Theorem}
\label{sec:main-proof}

\subsection{Sufficiency via fractional matching}

The first ingredient turns the setwise expansion in \Cref{eq:npdo} into a
fractional matching.  Its countable form requires both a finite approximation
and a tightness argument.

\begin{lemma}[Countable fractional matching lift]
\label{lem:matching}
Let $\Omega$ be finite or countable, let $p$ and $q$ be probability mass
functions on $\Omega$, let $R\subseteq\Omega\times\Omega$ be a relation, and
let $c\geq1$ and $\delta\in[0,1]$.  Suppose that
\begin{equation}
 p(S)\leq c q(R(S))+\delta
 \qquad\text{for every }S\subseteq\Omega.
 \label{eq:hall}
\end{equation}
Then there is a $c$-dominated fractional $R$-subcoupling $w$ of $p$ and $q$
such that
\begin{equation}
 \operatorname{mass}(w)\geq1-\delta.
 \label{eq:w-mass}
\end{equation}
\end{lemma}

\begin{proof}
Suppose first that $\Omega$ is finite.  Form a flow network with a source, a
left copy of $\Omega$, a right copy of $\Omega$, and a sink.  Give the edge
from the source to left vertex $a$ capacity $p(a)$, each middle edge from $a$
to $b$ with $aRb$ capacity $2+c$, and the edge from right vertex $b$ to the
sink capacity $cq(b)$.  The cut having only the source on its source side has
capacity $1$.  Hence a minimum cut cannot cross a middle edge, whose capacity
is greater than $1$.  If $S$ is the set of left vertices on the source side of
such a cut, all of $R(S)$ must also lie on its source side.  The cut therefore
has capacity at least
\[
 p(\Omega\setminus S)+cq(R(S))
 \geq 1-p(S)+p(S)-\delta
 =1-\delta,
\]
where \Cref{eq:hall} gives the inequality.
By the finite max-flow/min-cut fact in \Cref{fact:basic-tools}, there is a flow
of value at least $1-\delta$.  The flows on the middle edges define the
required subcoupling $w$.

Now suppose that $\Omega$ is countable.  By the summable-tail fact in
\Cref{fact:basic-tools}, choose finite sets
$F_n\subseteq\Omega$ and positive numbers $\rho_n\downarrow0$ such that
$p(F_n)\geq1-\rho_n$.  For each of the finitely many subsets
$S\subseteq F_n$, choose a finite set $G_{n,S}\subseteq R(S)$ satisfying
\[
 cq(R(S)\setminus G_{n,S})<\rho_n,
\]
and let $G_n$ be the union of these sets.  Apply the finite network construction
with left vertex set $F_n$ and right vertex set $G_n$.  A cut determined by
$S\subseteq F_n$ has capacity at least
\begin{align*}
 p(F_n\setminus S)+cq(R(S)\cap G_n)
 &\geq p(F_n)-p(S)+cq(R(S))-\rho_n\\
 &\geq p(F_n)-\delta-\rho_n\\
 &\geq1-\delta-2\rho_n.
\end{align*}
Consequently there is a function $w_n$, extended by zero to $\Omega^2$, that
satisfies \Cref{eq:w-support,eq:w-row,eq:w-column} and has total mass at least
$1-\delta-2\rho_n$.

Since $0\leq w_n(a,b)\leq p(a)\leq1$, the pointwise subsequence fact in
\Cref{fact:basic-tools}, applied to the countable set $\Omega^2$, gives a
subsequence along which every coordinate of $w_n$ converges.  Denote the
pointwise limit by $w$.  Passing to limits in finite partial sums proves
\Cref{eq:w-support,eq:w-row,eq:w-column}.
To retain the total mass, fix $\rho>0$ and choose finite
$F,G\subseteq\Omega$ such that
\[
 p(F^{\mathsf c})<\rho
 \qquad\text{and}\qquad
 cq(G^{\mathsf c})<\rho.
\]
The row and column bounds imply
\[
 \sum_{(a,b)\notin F\times G}w_n(a,b)
 \leq p(F^{\mathsf c})+cq(G^{\mathsf c})<2\rho.
\]
After passing to the limit in the finite sum over $F\times G$ and using
$\rho_n\to0$, we obtain
\[
 \sum_{a,b}w(a,b)\geq1-\delta-2\rho.
\]
Letting $\rho\downarrow0$ proves \Cref{eq:w-mass}.
\end{proof}

The matching lift is supported only on related executions.  It therefore
allows the admissibility inequality to be applied pointwise, while its row and
column bounds control the two execution laws and its missing mass accounts for
at most $\delta$.

\begin{lemma}[Contextual transport without error amplification]
\label{lem:transport}
Every $(\varepsilon,\delta)$-NPDO mechanism on a finite or countable execution
space satisfies, for every $\eta\geq0$, $\gamma\in[0,1]$, every
$(\eta,\gamma)$-admissible context $K$, every ordered pair
$x\sim_{\mathcal X}x'$, and every contextual event $A$,
\[
 (P_xK)(A)\leq e^{\varepsilon+\eta}(P_{x'}K)(A)
 +\delta+\gamma-\delta\gamma.
\]
In particular, it is $(\varepsilon,\delta)$-universally contextually
composable.
\end{lemma}

\begin{proof}
Fix an ordered pair $x\sim_{\mathcal X}x'$.  Apply \Cref{lem:matching} with
\[
 p=P_x,\qquad q=P_{x'},\qquad c=e^\varepsilon,
\]
and the execution relation $R$.  The hypothesis of that lemma is exactly
\Cref{eq:npdo}.  For the resulting $w$, set
\[
 r(a):=\sum_b w(a,b),\qquad
 s(b):=\sum_a w(a,b),\qquad
 m:=\sum_{a,b}w(a,b).
\]
Fix arbitrary $\eta\geq0$ and $\gamma\in[0,1]$, an observation space, an
$(\eta,\gamma)$-admissible kernel $K$, and an event $A$ in that space.  Write
$k(a):=K(a,A)$.  Since $0\leq k(a)\leq1$, the row and mass bounds give
\begin{align*}
 (P_xK)(A)
 &=\sum_a r(a)k(a)+\sum_a(P_x(a)-r(a))k(a)\\
 &\leq\sum_{a,b}w(a,b)k(a)+1-m.
\end{align*}
On the support of $w$, \Cref{eq:admissible} gives
$k(a)\leq e^\eta k(b)+\gamma$.  Thus, using the column bound and
$m\geq1-\delta$,
\begin{align*}
 (P_xK)(A)
 &\leq e^\eta\sum_b s(b)k(b)+\gamma m+1-m\\
 &\leq e^{\varepsilon+\eta}\sum_bP_{x'}(b)k(b)+1-(1-\gamma)m\\
 &\leq e^{\varepsilon+\eta}(P_{x'}K)(A)
       +1-(1-\gamma)(1-\delta)\\
 &=e^{\varepsilon+\eta}(P_{x'}K)(A)
       +\delta+\gamma-\delta\gamma.
\end{align*}
All choices were arbitrary.  Since $\delta+\gamma-\delta\gamma\leq
\delta+\gamma$, \Cref{eq:ucc} follows.
\end{proof}

It remains to prove that the implication is strict.  The next construction
places no mass at the middle vertex of a three-point path.  The endpoint laws
cannot satisfy a one-hop NPDO matching, but every admissible context relates
the endpoints through the middle vertex.

\subsection{The finite separating mechanism}

Let
\[
 \mathcal X=\{x_0,x_1\},\qquad
 \mathcal V=\{v\},\qquad
 \mathcal Y=\{0,1,2\}.
\]
Relate every pair of inputs, and define the output relation by
\[
 i\sim_{\mathcal Y}j
 \quad\Longleftrightarrow\quad
 |i-j|\leq1.
\]
Define a mechanism $M_\star$ by the execution laws
\begin{equation}
\begin{array}{c|ccc}
 & (v,0)&(v,1)&(v,2)\\ \hline
 P_{x_0}&1/3&0&2/3\\
 P_{x_1}&2/3&0&1/3
\end{array}.
\label{eq:three-point-law}
\end{equation}
Assume that an exact uniform sample from a three-element set is available at
unit cost.  Fix a constant-length access schedule $\sigma$, independent of the
input and sampled value, that writes to a designated output location, and take
$v$ to be its access-pattern view.  The following procedure implements the
already defined law.

\begin{algbox}
\phantomsection\label{alg:three-point}
\textbf{\uline{\textsc{ThreePoint}}}: Exact sampling of the endpoint-exchange mechanism.

\medskip
\noindent\uline{Input}: An input $x_i\in\mathcal X$, where $i\in\{0,1\}$.\\
\uline{Output}: An element $y\in\mathcal Y$ whose execution record is $(v,y)$.\\
\uline{Guarantee}: The execution law is the row indexed by $x_i$ in
\Cref{eq:three-point-law}.

\begin{enumerate}[leftmargin=*,itemsep=2pt,topsep=4pt]
  \item Sample $U$ exactly and uniformly from $\{1,2,3\}$.
  \item If $i=0$, assign $y\gets2$ for $U\leq2$ and $y\gets0$ for $U=3$.
        If $i=1$, assign $y\gets0$ for $U\leq2$ and $y\gets2$ for $U=3$.
  \item Execute $\sigma$ to write $y$ to the designated output location, and
        return $y$; the execution view is $v$.
\end{enumerate}
\end{algbox}

The implementation guarantee follows by counting the three equally likely
samples.  To prove the privacy claim, the essential additional step is to
convert the two adjacent admissibility inequalities into a two-hop endpoint
bound and then average according to \Cref{eq:three-point-law}.

\begin{lemma}[Exact realization and strict three-point separation]
\label{lem:separation}
The \hyperref[alg:three-point]{\textsc{ThreePoint}} algorithm realizes
\Cref{eq:three-point-law} in $O(1)$ time under exact uniform ternary sampling.
The resulting mechanism $M_\star$ is $(0,0)$-universally contextually
composable.  For every $\eta\geq0$ and $\gamma\in[0,1]$, each nontrivial
contextual comparison has additive term at most
$\max\{0,2+e^\eta-e^{2\eta}\}\gamma/3$, and hence at most $2\gamma/3$.
The mechanism is not $(0,0)$-NPDO.
\end{lemma}

\begin{proof}
For input $x_0$, two values of $U$ produce output $2$ and one produces output
$0$; for input $x_1$, these probabilities are exchanged.  The schedule
$\sigma$ always gives the fixed view $v$, so the joint laws are exactly
\Cref{eq:three-point-law}.  The procedure uses one unit-cost sample and a
constant number of other operations, hence runs in $O(1)$ time.

Fix $\eta\geq0$, $\gamma\in[0,1]$, an $(\eta,\gamma)$-admissible kernel $K$,
and an event $A$ in its observation space.  Set
\[
 t:=e^\eta,\qquad
 a:=K((v,0),A),\qquad
 b:=K((v,1),A),\qquad
 c:=K((v,2),A).
\]
The output relation is symmetric and contains the two adjacent pairs, so
admissibility in both directions yields
\begin{equation}
 c\leq tb+\gamma\leq t^2a+(t+1)\gamma,
 \qquad
 a\leq tb+\gamma\leq t^2c+(t+1)\gamma.
 \label{eq:two-hop}
\end{equation}
The probabilities of $A$ under the two contextual laws are
\[
 \mu:=(P_{x_0}K)(A)=\frac{a+2c}{3},
 \qquad
 \nu:=(P_{x_1}K)(A)=\frac{2a+c}{3}.
\]
Direct calculation gives
\begin{equation}
 3(\mu-t\nu)=(1-2t)a+(2-t)c.
 \label{eq:endpoint-difference}
\end{equation}
If $t\geq2$, both coefficients on the right-hand side are nonpositive.  If
$1\leq t<2$, the coefficient of $c$ is positive, so the first bound in
\Cref{eq:two-hop} gives
\begin{align*}
 3(\mu-t\nu)
 &\leq \bigl(1-2t+(2-t)t^2\bigr)a+(2-t)(t+1)\gamma\\
 &=- (t-1)(t^2-t+1)a+(2+t-t^2)\gamma\\
 &\leq(2+t-t^2)\gamma.
\end{align*}
Combining the two cases gives
\[
 \mu\leq t\nu+\frac{\max\{0,2+t-t^2\}}{3}\,\gamma.
\]
Exchanging $a$ and $c$ and using the second bound in \Cref{eq:two-hop} gives
the same inequality with $\mu$ and $\nu$ interchanged.  The inequalities for
the two reflexive input pairs hold with zero additive term because $t\geq1$.
For $t\geq1$, the displayed coefficient is at most $2/3$; since
$2\gamma/3\leq\gamma$, this proves $(0,0)$-universal contextual
composability, as well as the stated sharper comparison.

Finally, let $S=\{(v,2)\}$.  Its relational neighborhood is
\[
 R(S)=\{(v,1),(v,2)\}.
\]
\Cref{eq:three-point-law} gives
\[
 P_{x_0}(S)=\frac23>\frac13=P_{x_1}(R(S)),
\]
which contradicts \Cref{eq:npdo} with $\varepsilon=\delta=0$.  Thus
$M_\star$ is not $(0,0)$-NPDO.
\end{proof}

\Cref{lem:transport} proves the sufficiency assertion of \Cref{thm:main}, and
\Cref{lem:separation} supplies the finite constant-time counterexample.  Hence
universal contextual composability is strictly weaker than NPDO in the stated
finite-or-countable statistical execution model.

\bibliographystyle{alpha}
\bibliography{references}

\end{document}
